% =====================================================================
%  Chapter 4: The Linear Model
%  Source: ClassNormReg.tex (Overleaf) + P3 transcription
% =====================================================================

\chapter[The Linear Model]{The Classical Linear Model}
\label{chap_LinearModel}

\section{The Regression Function}
\label{sec_RegFunction}

The general regression model is:
\begin{equation}
  y = m(\mathbf{x}) + e, \qquad
  m(\mathbf{x}) \equiv \E[y \mid \mathbf{x}]
  = \int_{-\infty}^{\infty} y\, f(y \mid \mathbf{x})\, dy
\end{equation}

By construction, the regression error $e = y - m(\mathbf{x})$ satisfies:
\begin{enumerate}[label=(\roman*)]
  \item $\E[e \mid \mathbf{x}] = 0$
  \item $\E[e] = 0$
  \item $\E[e\mathbf{x}] = 0$
\end{enumerate}

\subsection{Homoskedasticity vs.\ Heteroskedasticity}

The term \emph{skedasis} means dispersion. Two cases arise for the conditional
error variance $\E[e^2 \mid \mathbf{x}]$:

\begin{enumerate}[label=\alph*)]
  \item \textbf{Homoskedasticity:}\; $\E[e^2\mid\mathbf{x}] = \sigma^2$
    (constant conditional variance)
  \item \textbf{Heteroskedasticity:}\;
    $\E[e^2\mid\mathbf{x}] = \sigma^2(\mathbf{x})$
    (skedastic function; variance depends on $\mathbf{x}$)
\end{enumerate}

The unconditional variance via the law of iterated expectations:
\begin{equation}
  \sigma^2 = \E[e^2] = \E\!\bigl[\E[e^2 \mid \mathbf{x}]\bigr]
  = \E[\sigma^2(\mathbf{x})]
\end{equation}

The \textbf{linear model}\index{linear model} specifies:
\begin{alignat*}{2}
  y &= \mathbf{x}'\beta + e \\
  \E[e \mid \mathbf{x}] &= 0 \\
  \E[e^2 \mid \mathbf{x}] &= \sigma^2 \quad\text{or}\quad \sigma^2(\mathbf{x})
\end{alignat*}

\noindent\textbf{Modelling conditional variance.}\index{ARCH}\index{GARCH}
Robert Engle's ARCH model (1982), for which he received the 2003 Nobel Prize,
specifies $\E[e^2 \mid \mathbf{x}] = \sigma^2(\mathbf{x})$ as an explicit
function of past squared errors. The generalised GARCH model extends this to
allow the conditional variance to depend on past conditional variances as well.
These models are central to financial time series but also arise in spatial
panels where variance may cluster geographically.


\section{The Best Linear Predictor}
\label{sec_BLP}

Minimise the expected squared prediction error over $\beta$:
\begin{equation}
  F(\beta) = \E\!\bigl[(y - \mathbf{x}'\beta)^2\bigr]
  = \E[y^2] - 2\beta'\E[\mathbf{x}y] + \beta'\E[\mathbf{x}\mathbf{x}']\beta
\end{equation}

\noindent\textbf{First order condition:}
$\partial F / \partial\beta = -2\E[\mathbf{x}y] + 2\E[\mathbf{x}\mathbf{x}']\beta = 0$.

If $\E[\mathbf{x}\mathbf{x}']$ is invertible:
\begin{tcolorbox}[keyresult]
\begin{equation}
  \beta = \bigl(\E[\mathbf{x}\mathbf{x}']\bigr)^{-1}\E[\mathbf{x}y]
\end{equation}
\end{tcolorbox}

\subsection{The Empirical Analogue: OLS}

Replace population expectations with sample averages:
\begin{equation}
  F_n(\beta) = \frac{1}{n}\sum(y - \mathbf{x}'\beta)^2
  = \frac{1}{n}\,\text{SSE}_n(\beta)
\end{equation}

Setting $\partial\,\text{SSE}/\partial\beta = 0$:
\begin{equation}
  \hat\beta = \Bigl(\sum_i\mathbf{x}_i\mathbf{x}_i'\Bigr)^{-1}
              \Bigl(\sum_i\mathbf{x}_i y_i\Bigr)
\end{equation}


\section{The Model in Matrix Notation}
\label{sec_MatrixNotation}

The system of $n$ equations $y_i = \mathbf{x}_i'\beta + e_i$ is written
compactly as:
\begin{equation}
  \underset{n\times1}{\mathbf{y}}
  = \underset{n\times k}{X}\;\underset{k\times1}{\beta}
  + \underset{n\times1}{\mathbf{e}}
\end{equation}

\begin{tcolorbox}[keyresult]
\textbf{OLS estimator (matrix form):}
\begin{equation}
  \hat\beta = (X'X)^{-1}X'\mathbf{y}
  \label{eq_OLS}
\end{equation}
Residuals: $\hat{\mathbf{e}} = \mathbf{y} - X\hat\beta$.\quad
Variance estimate: $\hat\sigma^2 = n^{-1}\hat{\mathbf{e}}'\hat{\mathbf{e}}$.
\end{tcolorbox}

\noindent\textbf{Three ways to write the linear model:}
\begin{enumerate}
  \item Equation by equation:\; $y_i = \mathbf{x}_i'\beta + e_i$
  \item Matrix form:\; $\mathbf{y} = X\beta + \mathbf{e}$
  \item Scalar form:\; $y_i = \beta_0 + \beta_1 x_{1i} + \cdots
    + \beta_{k-1}x_{k-1,i} + e_i$
\end{enumerate}


\section{Sampling Distribution of $\hat\beta$}
\label{sec_SamplingDist}

$\hat\beta$ is a random variable with its own $\mu_{\hat\beta}$ and
$\sigma^2_{\hat\beta}$, and therefore has a \textbf{sampling
distribution}\index{sampling distribution}.

\subsection{Unbiasedness of OLS}

\begin{equation}
  \E(\hat\beta \mid X) = \beta \qquad\text{and}\qquad \E(\hat\beta) = \beta
\end{equation}

\noindent\textit{Proof:}
\begin{align*}
  \E(\hat\beta \mid X)
  &= \E\!\bigl[(X'X)^{-1}X'\mathbf{y} \mid X\bigr] \\
  &= (X'X)^{-1}X'\,\E[\mathbf{y} \mid X]
   = (X'X)^{-1}X'X\beta = \beta \qquad\square
\end{align*}

\subsection{Variance of $\hat\beta$}

Under homoskedasticity, $\E[\mathbf{e}\mathbf{e}'] = \sigma^2 I_n$:

\begin{tcolorbox}[keyresult]
\begin{equation}
  V(\hat\beta) = \sigma^2(X'X)^{-1}
  \label{eq_VarBeta}
\end{equation}
\end{tcolorbox}

\noindent\textit{Proof:} $\hat\beta - \beta = (X'X)^{-1}X'\mathbf{e}$, so:
\begin{align*}
  V(\hat\beta)
  &= \E\!\bigl[(\hat\beta-\beta)(\hat\beta-\beta)'\bigr]
   = (X'X)^{-1}X'\,\E[\mathbf{e}\mathbf{e}']\,X(X'X)^{-1} \\
  &= (X'X)^{-1}X'\,\sigma^2 I\,X(X'X)^{-1}
   = \sigma^2(X'X)^{-1} \qquad\square
\end{align*}

\subsection{Estimating $\sigma^2$}

OLS does not directly provide an estimate of $\sigma^2$. Using residuals
$\hat{e}_i$ in place of errors $e_i$ yields a downward-biased estimate:
\begin{equation}
  \E[\hat\sigma^2] = \frac{n-k}{n}\,\sigma^2 < \sigma^2
\end{equation}

\begin{tcolorbox}[keyresult]
\begin{equation}
  s^2 \equiv \frac{1}{n-k}\sum_{i=1}^{n}\hat{e}_i^2
  = \frac{n}{n-k}\,\hat\sigma^2
  \qquad \text{(unbiased estimator of }\sigma^2\text{)}
\end{equation}
\end{tcolorbox}

Note: while $s^2$ is unbiased for $\sigma^2$, $s$ is \emph{not} unbiased
for $\sigma$.


\section{Geometry of OLS}
\label{sec_GeometryOLS}

Expanding $\hat\beta = (X'X)^{-1}X'\mathbf{y}$:
\begin{align*}
  \hat\beta
  &= (X'X)^{-1}X'(X\beta + \mathbf{e}) \\
  &= \beta + (X'X)^{-1}X'\mathbf{e}
\end{align*}
Fitted values: $\hat{\mathbf{y}} = X\hat\beta = X(X'X)^{-1}X'\mathbf{y} =
P_X\mathbf{y}$.

\subsection{Projection Matrices and Their Relationships}
\label{subsec_Projections}

The \textbf{hat matrix}\index{hat matrix} and its complement
(Section~\ref{sec_Projections}) satisfy additional structural relationships
when $X = [X_1\;X_2]$. Defining $P_1 = X_1(X_1'X_1)^{-1}X_1'$:
\begin{align}
  P_X P_1 &= P_1 \qquad\text{since}\quad P_1 < P_X \label{eq_PP1} \\
  M_X M_1 &= M_X \qquad\text{since}\quad M_1 > M_X \label{eq_MM1}
\end{align}

The orthogonal decomposition $\mathbf{y} = P_X\mathbf{y} + M_X\mathbf{y}
= \hat{\mathbf{y}} + \hat{\mathbf{e}}$ is illustrated in
Figure~\ref{fig_GeomOLS}.

\begin{figure}[htp]
\centering
\begin{tikzpicture}[>=Stealth, scale=1.3]
  \fill[gray!12]
    (-0.5,-0.35) -- (3.7,-0.35) -- (4.2,0.95) -- (0.0,0.95) -- cycle;
  \node[gray!70, font=\small] at (1.8, -0.05) {column space of $X$};
  \coordinate (O)   at (1.4, 0.15);
  \fill (O) circle (1.5pt) node[below left] {$O$};
  \coordinate (Y)   at (3.4, 3.1);
  \coordinate (PXY) at (3.4, 0.60);
  \draw[->, very thick]     (O) -- (Y)   node[right] {$\mathbf{y}$};
  \draw[->, thick]          (O) -- (PXY) node[below right]
    {$\hat{\mathbf{y}} = P_X\mathbf{y}$};
  \draw[->, thick, dashed] (PXY) -- (Y)  node[midway, right=2pt]
    {$\hat{\mathbf{e}} = M_X\mathbf{y}$};
  \draw ($(PXY)+(-0.17,0)$) -- ++(0,0.17) -- ++(0.17,0);
  \coordinate (X1) at (2.8, 0.42);
  \coordinate (X2) at (1.9, 0.80);
  \draw[->, thick] (O) -- (X1) node[below] {$X_1$};
  \draw[->, thick] (O) -- (X2) node[above left] {$X_2$};
\end{tikzpicture}
\caption{Geometry of OLS. $\hat{\mathbf{y}} = P_X\mathbf{y}$ is the
orthogonal projection of $\mathbf{y}$ onto the column space of $X$; the
residual vector $\hat{\mathbf{e}} = M_X\mathbf{y}$ is perpendicular to
that space. Note $\hat{\mathbf{y}}'\hat{\mathbf{e}} = 0$.}
\label{fig_GeomOLS}
\end{figure}


\section{Goodness of Fit}
\label{sec_GoodnessOfFit}

\subsection{The TSS Decomposition}

\begin{tcolorbox}[keyresult]
\begin{alignat*}{2}
  \text{TSS} &= \mathbf{y}'M_\iota\mathbf{y}
    = \textstyle\sum(y_i-\bar{y})^2 \\
  \text{ESS} &= \mathbf{y}'M_\iota P_X\mathbf{y}
    = \textstyle\sum(\hat{y}_i-\bar{y})^2 \\
  \text{RSS} &= \mathbf{y}'M_X\mathbf{y}
    = \textstyle\sum(y_i-\hat{y}_i)^2
\end{alignat*}
\vspace{-4pt}
\begin{equation}
  R^2 = \frac{\text{ESS}}{\text{TSS}} = 1 - \frac{\text{RSS}}{\text{TSS}}
\end{equation}
\end{tcolorbox}

\subsection{Adjusted $R^2$}

\begin{equation}
  \bar{R}^2 = 1 - (1-R^2)\frac{n-1}{n-k}
\end{equation}

$\bar{R}^2$ increases only if the $t$-statistic of the newly added variable
exceeds $|1|$ in absolute value. An increase in $R^2$ is \emph{not}
indicative of improved specification---even poorly specified equations can
achieve high $R^2$.


\section{The Frisch--Waugh--Lovell Theorem}
\label{sec_FWL}

\begin{tcolorbox}[definition]
\textbf{Theorem (FWL; see \citealt{Lovell2008,Davidson2004}).} If
$Y = X_1\beta_1 + X_2\beta_2 + \mathbf{e}$ ($X_1$ is $n\times k_1$;
$X_2$ is $n\times k_2$), then the OLS estimate of $\beta_2$ from the full
regression equals the OLS estimate from the auxiliary regression:
\begin{equation}
  M_{X_1}Y = M_{X_1}X_2\beta_2 + \mathbf{v}, \quad \mathbf{v} = M_{X_1}\mathbf{e}
\end{equation}
where $M_{X_1} = I - X_1(X_1'X_1)^{-1}X_1'$.
\end{tcolorbox}

\noindent\textbf{Proof.} Write $M_{X_1} \equiv M_1$.
Premultiply $y = X_1\hat\beta_1 + X_2\hat\beta_2 + M_Xy$ by $X_2'M_1$:
\begin{equation}
  X_2'M_1y
  = \underbrace{X_2'M_1X_1\hat\beta_1}_{(1)}
  + \underbrace{X_2'M_1X_2\hat\beta_2}_{(2)}
  + \underbrace{X_2'M_1M_Xy}_{(3)}
\end{equation}
Two properties eliminate terms (1) and (3):
\begin{enumerate}
  \item $M_1X_1 = 0$ $\Rightarrow$ term (1) $= 0$.
  \item $M_1M_X = M_X$ (since $M_X < M_1$) $\Rightarrow$
    $X_2'M_1M_Xy = X_2'M_Xy = X_2'\hat{\mathbf{e}} = 0$.
\end{enumerate}

\begin{tcolorbox}[keyresult]
\begin{equation}
  \hat\beta_2 = (X_2'M_1X_2)^{-1}X_2'M_1y,
  \qquad
  \Var(\hat\beta_2) = \sigma^2(X_2'M_1X_2)^{-1}
  \label{eq_FWL}
\end{equation}
\end{tcolorbox}

\noindent\textbf{Interpretation.} The FWL theorem says that the coefficient
on $X_2$ in the full regression is numerically identical to the coefficient
from regressing the residuals of $Y$ on $X_1$ onto the residuals of $X_2$
on $X_1$. In other words, OLS \emph{partials out} all other regressors before
estimating any particular coefficient. A regression coefficient does not
measure the raw correlation between $Y$ and $X_2$; it measures the
correlation after removing everything both $Y$ and $X_2$ share with $X_1$.


\section{Specification Issues}
\label{sec_Specification}

\subsection{Overspecification}

Variables not belonging to the true DGP are mistakenly included.
Overspecification is \textbf{not} misspecification.

Suppose we estimate $y = X\beta + Z\gamma + \mathbf{e}$,
$\mathbf{e}\sim\text{iid}(0,\sigma^2I)$, but the true DGP is
$y = X\beta_0 + \mathbf{e}$. The overspecified model nests the true model
(setting $\gamma=0$), so $\hat\beta$ is unbiased.

\begin{tcolorbox}[keyresult]
\textbf{Overspecification:} unbiased, (weakly) consistent, nominally
inefficient estimates.
\end{tcolorbox}

\subsection{Underspecification and Omitted Variable Bias}
\label{subsec_OVB}

We estimate $y = X\beta + \mathbf{e}$ but the true DGP is:
\begin{equation}
  y = X\beta_0 + Z\gamma_0 + \mathbf{e},
  \quad \mathbf{e}\sim\text{iid}(0,\sigma_0^2I)
\end{equation}

\noindent\textbf{Omitted variable bias:}\index{omitted variable bias}
\begin{align}
  \E(\hat\beta)
  &= \beta_0 + \underbrace{(X'X)^{-1}X'Z\gamma_0}_{\text{OVB}}
\end{align}

\begin{tcolorbox}[keyresult]
The OVB is zero if and only if:
\begin{enumerate}
  \item $X'Z = 0$ (omitted variable is orthogonal to included regressors), or
  \item $\gamma_0 = 0$ (the omitted variable has no effect on $y$).
\end{enumerate}
\end{tcolorbox}

\begin{tcolorbox}[keyresult]
\textbf{Underspecification:} biased, consistent, linear estimator.
\end{tcolorbox}


\section{Functional Form and Nonlinearity}
\label{sec_FunctionalForm}

The linear model requires linearity in \emph{parameters}, not in variables.
This distinction permits substantial flexibility.

\subsection{Polynomial Terms}

\begin{equation}
  y = \beta_0 + \beta_1 x + \beta_2 x^2 + e
\end{equation}

Substituting $z = x^2$ preserves linearity in the parameters while capturing
a nonlinear relationship in $x$.

\subsection{Log-linearisation: The Cobb-Douglas Production Function}

The \textbf{Cobb-Douglas}\index{Cobb-Douglas production function} function:
\begin{equation}
  y = AK^\alpha L^\beta e, \qquad\text{with CES: }\alpha = 1-\beta
\end{equation}

Log-linearisation yields a standard linear regression model:
\begin{equation}
  \ln y = \ln A + (1-\beta)\ln K + \beta\ln L + \ln e
\end{equation}

\noindent\textbf{Constant returns to scale} ($F(aK,aL) = aF(K,L)$):
\begin{equation}
  F(aK,aL) = A(aK)^\alpha(aL)^{1-\alpha} = a\,AK^\alpha L^{1-\alpha}
  \;\checkmark
\end{equation}
If $\alpha + \beta > 1$: increasing returns to scale.

\noindent\textbf{Elasticity:}
\begin{equation}
  \frac{dy/dL}{y/L}
  = \frac{\text{marginal product}}{\text{average product}}
  = \beta \quad\text{(constant factor proportion)}
\end{equation}

\noindent\textbf{Derivation from CES.} The Cobb-Douglas arises as a limiting
case of the CES function $Y = A[\alpha K^\rho + (1-\alpha)L^\rho]^{\gamma/\rho}$
as $\rho \to 0$. Applying L'H\^opital's rule to $\ln Y$:
\begin{equation}
  \lim_{\rho\to0}\ln Y = \ln A + \alpha\ln K + (1-\alpha)\ln L
  \;\Longrightarrow\; Y = AK^\alpha L^{1-\alpha}
\end{equation}


\section{Maximum Likelihood Estimation and OLS}
\label{sec_MLE}

The \textbf{classical normal regression} (CNR) model adds the assumption that
$e_i \mid \mathbf{x}_i \sim \mathcal{N}(0, \sigma^2)$, implying:
\begin{equation}
  y_i \mid \mathbf{x}_i \sim \mathcal{N}(\mathbf{x}_i'\beta, \sigma^2)
\end{equation}

The likelihood function for observation $i$:
\begin{equation}
  \mathcal{L}_i = f(y_i) = \frac{1}{\sqrt{2\pi\sigma^2}}
  \exp\!\left(-\frac{(y_i - \mathbf{x}_i'\beta)^2}{2\sigma^2}\right)
\end{equation}

The \textbf{maximum likelihood estimator}\index{maximum likelihood
estimator (MLE)} $(\hat\beta_{mle}, \hat\sigma^2_{mle})$ maximises
$\log\mathcal{L}(\beta,\sigma^2)$. Since this is a function of $\beta$ only
through $\text{SSE}_n(\beta)$, maximising the likelihood is equivalent to
minimising the sum of squared errors:

\begin{tcolorbox}[keyresult]
\begin{equation}
  \hat\beta_{mle} = \hat\beta_{ols}
\end{equation}
\end{tcolorbox}

Under the normality assumption, OLS is not merely the best linear unbiased
estimator (BLUE by the Gauss-Markov theorem); it is the minimum variance
unbiased estimator among \emph{all} unbiased estimators.


\section{General Method of Moments}
\label{sec_GMM}

OLS exploits the moment condition $\E[\mathbf{x}_i e_i] = 0$.
The \textbf{General Method of Moments}\index{GMM|see{General Method of
Moments}}\index{General Method of Moments} (GMM) generalises this by
allowing for a vector of moment conditions $\E[\mathbf{z}_i g(\mathbf{y}_i,
\boldsymbol{\theta})] = \mathbf{0}$, where $\mathbf{z}_i$ may be instruments
distinct from regressors $\mathbf{x}_i$. OLS is a special case of GMM with
$\mathbf{z}_i = \mathbf{x}_i$; IV/2SLS is another special case with
$\mathbf{z}_i$ containing external instruments (Chapter~\ref{chap_BeyondOLS}).

% TODO: Expand GMM with J-test for over-identification

\renewcommand\bibname{Further reading}
\begin{subappendices}
\bibliographystyle{econometrica}
\bibliography{Literature/OverLord}
\IfFileExists{Literature/ClassNormReg.bbl}{\chapter[The CNR model]{The classical normal regression (CNR) model}\label{chap_ClassNormReg}

\section{The CNR model}
\subsection{Basic set-up}
\subsection{The model in matrix notation}

\section{Different approaches to OLS}
\subsection{Least squares estimator}

\subsection{Geometric approach to OLS}
\subsubsection{Projection matrices and their relationships}
Let the matrix $\boldsymbol{P_{X}}$ be defined as

\begin{eqnarray}
\boldsymbol{P_{X}} \equiv \boldsymbol{X}(\boldsymbol{X}'\boldsymbol{X})^{-1}\boldsymbol{X}'.\label{eqn_CNR_Px}
\end{eqnarray}

$\boldsymbol{P_{X}}$ belongs to the general class of so-called projection matrices which have a number of interesting properties. Most importantly, the projection matrix $\boldsymbol{P_{X}}$ is a square matrix of dimension $n \times n$ that gives a vector space projection from $\mathbb{R}^{n}$ to a given subspace. Projection matrices are said to be \emph{idempotent} which implies that $\boldsymbol{P_{X}}\boldsymbol{P_{X}}=\boldsymbol{P_{X}}$.

This important property is straightforward to demonstrate in the case of equation \eqref{eqn_CNR_Px}, where we can use the fact that $$(\boldsymbol{X}'\boldsymbol{X})^{-1}\boldsymbol{X}'\boldsymbol{X}=\boldsymbol{I}$$ in order to show that $$\boldsymbol{P_{X}}\boldsymbol{P_{X}}=\boldsymbol{X}(\boldsymbol{X}'\boldsymbol{X})^{-1}\boldsymbol{X}'\boldsymbol{X}(\boldsymbol{X}'
\boldsymbol{X})^{-1}\boldsymbol{X}'=\boldsymbol{X}(\boldsymbol{X}'\boldsymbol{X})^{-1}\boldsymbol{X}'=\boldsymbol{P_{X}}$$. Furthermore, projection matrices can generally be classified either as \emph{orthogonal} projection matrices or \emph{oblique} (non-orthogonal) projection matrices.

In addition to being square, a projection matrix is also a symmetric matrix if and only if the vector space projection is orthogonal. $\boldsymbol{P_{X}}$ is an orthogonal projection matrix that spans the space $\mathcal{S}(\boldsymbol{X})$; consequently, any point in $\boldsymbol{X}$ that is projected by $\boldsymbol{P_{X}}$ will be projected onto itself. This implies that $\boldsymbol{P_{X}}\boldsymbol{X} = \boldsymbol{X}$. Again, this is easily verified by using equation \eqref{eqn_CNR_Px} which yields $\boldsymbol{P_{X}}\boldsymbol{X}=\boldsymbol{X}(\boldsymbol{X}'\boldsymbol{X})^{-1}\boldsymbol{X}'\boldsymbol{X} = \boldsymbol{X}$.

$\boldsymbol{M_{X}}$ is the orthogonal complement to $\boldsymbol{P_{X}}$, such that $\boldsymbol{M_{X}} = (\boldsymbol{I} - \boldsymbol{P_{X}})$. Consequently,

\begin{eqnarray}
\boldsymbol{P_{X}}\boldsymbol{M_{X}} &=& \boldsymbol{P_{X}}(\boldsymbol{I} - \boldsymbol{P_{X}}) \nonumber \\
&=& \boldsymbol{P_{X}}-\boldsymbol{P_{X}}\boldsymbol{P_{X}} = \boldsymbol{P_{X}}-\boldsymbol{P_{X}}=\boldsymbol{0}.
\end{eqnarray}

Since $\boldsymbol{P_{X}}$ and $\boldsymbol{M_{X}}$ are orthogonal to each other, any vector in $\mathcal{S}(\boldsymbol{X})$ will either be projected onto itself (i.e. $\boldsymbol{P_{X}}\boldsymbol{X} = \boldsymbol{X}$) or it will be annihilated ($\boldsymbol{M_{X}}\boldsymbol{X} = \boldsymbol{0}$).

Now, suppose that the matrix $\boldsymbol{X}$ can be decomposed into the submatrices $\boldsymbol{X}_{1}$ and $\boldsymbol{X}_{2}$, such that $\boldsymbol{X} = [\boldsymbol{X}_{1}\hspace{1ex}\boldsymbol{X}_{2}]$. Analogous to $\boldsymbol{P_{X}}$, we can now define the projection matrix $\boldsymbol{P}_{1} = \boldsymbol{X}_{1}(\boldsymbol{X}_{1}'\boldsymbol{X}_{1})^{-1}\boldsymbol{X}_{1}'$. This yields the following important relationships:

\begin{eqnarray}
\boldsymbol{P_{X}}\boldsymbol{P}_{1}=\boldsymbol{P}_{1} \qquad &\text{since}& \qquad \boldsymbol{P}_{1}<\boldsymbol{P_{X}}\\
\boldsymbol{M_{X}}\boldsymbol{M}_{1}=\boldsymbol{M_{X}} \qquad &\text{since}& \qquad \boldsymbol{M}_{1}>\boldsymbol{M_{X}}
\end{eqnarray}


\subsubsection{Geometry of the Frisch-Waugh-Lovell theorem}
See \citet{Lovell2008} for a simple proof that does not rely on matrix algebra. More comprehensive treatment of the Frisch-Waugh-Lovell theorem (FWL) can be found in \citet{Davidson2004}.

\subsection{General Method of Moments (GMM)}

\subsection{Maximum Likelihood Estimation and OLS}
The maximum likelihood estimate (MLE) is that assumed value of an unknown population parameter which maximises the probability of obtaining the sample data that is actually observed.

The normal regression model is the linear regression model under the restriction that the error term $e_{i}$ is independent of $\boldsymbol{x}_{i}$ and has the distribution $(0, \sigma^{2})$ Thus the standard set-up of OLS that we have just seen can easily be re-written as

\begin{eqnarray}
e_{i} | \boldsymbol{x}_{i} \sim (0, \sigma^{2})
\end{eqnarray}

Recall from section~\ref{chap_Review:Stats} that the mean and variance of a variable $V = \mu_{0} + e$ where $e \sim (0,\sigma^{2})$, is $\mu_{V} = \mu$ and $\sigma^{2}_{V}=\sigma^{2}_{e}$. This implies that

\begin{eqnarray}
y_{i} | \boldsymbol{x}_{i} \sim (\boldsymbol{x}_{i}'\boldsymbol{\beta}, \sigma^{2})
\end{eqnarray}

Equivalently, we can rewrite the linear model as $\boldsymbol{y} \sim N{\boldsymbol{\Omega}, \boldsymbol{\Omega}}$ where $\boldsymbol{\Omega}$ is assumed to meet the Gauss-Markov assumptions.

Given the parametric nature of the CNR, we can use likelihood methods\index{likelihood methods} for estimating the parameters. The likelihood function for each observation is

\begin{eqnarray}
\mathcal{L}_{i}  = f(y_{i}) = \frac{1}{\sqrt{2 \pi \sigma^{2}}}e^{-\frac{(y_{i} - \boldsymbol{x}'_{i}\boldsymbol{\beta})^{2}}{2\sigma^{2}}}
\end{eqnarray}

The maximum likelihood estimator\index{maximum likelihood estimator (MLE)|see{likelihood methods}} (MLE; $\hat{\boldsymbol{\beta}}_{mle}, \sigma^{2}_{mle}$) maximises $\log \mathcal{L}(\boldsymbol{\beta}, \sigma^{2})$. Since this is a function of $\boldsymbol{\beta}$ only through the sum of squared errors $SSE_{n}(\boldsymbol{\beta})$, maximising the likelihood is equivalent to minimising the sum of squared errors. Therefore

\begin{eqnarray}
\hat{\boldsymbol{\beta}}_{mle} = \boldsymbol{\hat{\beta}}_{ols}.
\end{eqnarray}


\begin{subappendices}
\section{Code for Spatial Stats in R}
\lstinputlisting{"Programs/UP504_OLS.R"}

\renewcommand\bibname{Further reading}

\bibliographystyle{econometrica}
\bibliography{Literature/OverLord}
\IfFileExists{Literature/ClassNormReg.bbl}{\chapter[The CNR model]{The classical normal regression (CNR) model}\label{chap_ClassNormReg}

\section{The CNR model}
\subsection{Basic set-up}
\subsection{The model in matrix notation}

\section{Different approaches to OLS}
\subsection{Least squares estimator}

\subsection{Geometric approach to OLS}
\subsubsection{Projection matrices and their relationships}
Let the matrix $\boldsymbol{P_{X}}$ be defined as

\begin{eqnarray}
\boldsymbol{P_{X}} \equiv \boldsymbol{X}(\boldsymbol{X}'\boldsymbol{X})^{-1}\boldsymbol{X}'.\label{eqn_CNR_Px}
\end{eqnarray}

$\boldsymbol{P_{X}}$ belongs to the general class of so-called projection matrices which have a number of interesting properties. Most importantly, the projection matrix $\boldsymbol{P_{X}}$ is a square matrix of dimension $n \times n$ that gives a vector space projection from $\mathbb{R}^{n}$ to a given subspace. Projection matrices are said to be \emph{idempotent} which implies that $\boldsymbol{P_{X}}\boldsymbol{P_{X}}=\boldsymbol{P_{X}}$.

This important property is straightforward to demonstrate in the case of equation \eqref{eqn_CNR_Px}, where we can use the fact that $$(\boldsymbol{X}'\boldsymbol{X})^{-1}\boldsymbol{X}'\boldsymbol{X}=\boldsymbol{I}$$ in order to show that $$\boldsymbol{P_{X}}\boldsymbol{P_{X}}=\boldsymbol{X}(\boldsymbol{X}'\boldsymbol{X})^{-1}\boldsymbol{X}'\boldsymbol{X}(\boldsymbol{X}'
\boldsymbol{X})^{-1}\boldsymbol{X}'=\boldsymbol{X}(\boldsymbol{X}'\boldsymbol{X})^{-1}\boldsymbol{X}'=\boldsymbol{P_{X}}$$. Furthermore, projection matrices can generally be classified either as \emph{orthogonal} projection matrices or \emph{oblique} (non-orthogonal) projection matrices.

In addition to being square, a projection matrix is also a symmetric matrix if and only if the vector space projection is orthogonal. $\boldsymbol{P_{X}}$ is an orthogonal projection matrix that spans the space $\mathcal{S}(\boldsymbol{X})$; consequently, any point in $\boldsymbol{X}$ that is projected by $\boldsymbol{P_{X}}$ will be projected onto itself. This implies that $\boldsymbol{P_{X}}\boldsymbol{X} = \boldsymbol{X}$. Again, this is easily verified by using equation \eqref{eqn_CNR_Px} which yields $\boldsymbol{P_{X}}\boldsymbol{X}=\boldsymbol{X}(\boldsymbol{X}'\boldsymbol{X})^{-1}\boldsymbol{X}'\boldsymbol{X} = \boldsymbol{X}$.

$\boldsymbol{M_{X}}$ is the orthogonal complement to $\boldsymbol{P_{X}}$, such that $\boldsymbol{M_{X}} = (\boldsymbol{I} - \boldsymbol{P_{X}})$. Consequently,

\begin{eqnarray}
\boldsymbol{P_{X}}\boldsymbol{M_{X}} &=& \boldsymbol{P_{X}}(\boldsymbol{I} - \boldsymbol{P_{X}}) \nonumber \\
&=& \boldsymbol{P_{X}}-\boldsymbol{P_{X}}\boldsymbol{P_{X}} = \boldsymbol{P_{X}}-\boldsymbol{P_{X}}=\boldsymbol{0}.
\end{eqnarray}

Since $\boldsymbol{P_{X}}$ and $\boldsymbol{M_{X}}$ are orthogonal to each other, any vector in $\mathcal{S}(\boldsymbol{X})$ will either be projected onto itself (i.e. $\boldsymbol{P_{X}}\boldsymbol{X} = \boldsymbol{X}$) or it will be annihilated ($\boldsymbol{M_{X}}\boldsymbol{X} = \boldsymbol{0}$).

Now, suppose that the matrix $\boldsymbol{X}$ can be decomposed into the submatrices $\boldsymbol{X}_{1}$ and $\boldsymbol{X}_{2}$, such that $\boldsymbol{X} = [\boldsymbol{X}_{1}\hspace{1ex}\boldsymbol{X}_{2}]$. Analogous to $\boldsymbol{P_{X}}$, we can now define the projection matrix $\boldsymbol{P}_{1} = \boldsymbol{X}_{1}(\boldsymbol{X}_{1}'\boldsymbol{X}_{1})^{-1}\boldsymbol{X}_{1}'$. This yields the following important relationships:

\begin{eqnarray}
\boldsymbol{P_{X}}\boldsymbol{P}_{1}=\boldsymbol{P}_{1} \qquad &\text{since}& \qquad \boldsymbol{P}_{1}<\boldsymbol{P_{X}}\\
\boldsymbol{M_{X}}\boldsymbol{M}_{1}=\boldsymbol{M_{X}} \qquad &\text{since}& \qquad \boldsymbol{M}_{1}>\boldsymbol{M_{X}}
\end{eqnarray}


\subsubsection{Geometry of the Frisch-Waugh-Lovell theorem}
See \citet{Lovell2008} for a simple proof that does not rely on matrix algebra. More comprehensive treatment of the Frisch-Waugh-Lovell theorem (FWL) can be found in \citet{Davidson2004}.

\subsection{General Method of Moments (GMM)}

\subsection{Maximum Likelihood Estimation and OLS}
The maximum likelihood estimate (MLE) is that assumed value of an unknown population parameter which maximises the probability of obtaining the sample data that is actually observed.

The normal regression model is the linear regression model under the restriction that the error term $e_{i}$ is independent of $\boldsymbol{x}_{i}$ and has the distribution $(0, \sigma^{2})$ Thus the standard set-up of OLS that we have just seen can easily be re-written as

\begin{eqnarray}
e_{i} | \boldsymbol{x}_{i} \sim (0, \sigma^{2})
\end{eqnarray}

Recall from section~\ref{chap_Review:Stats} that the mean and variance of a variable $V = \mu_{0} + e$ where $e \sim (0,\sigma^{2})$, is $\mu_{V} = \mu$ and $\sigma^{2}_{V}=\sigma^{2}_{e}$. This implies that

\begin{eqnarray}
y_{i} | \boldsymbol{x}_{i} \sim (\boldsymbol{x}_{i}'\boldsymbol{\beta}, \sigma^{2})
\end{eqnarray}

Equivalently, we can rewrite the linear model as $\boldsymbol{y} \sim N{\boldsymbol{\Omega}, \boldsymbol{\Omega}}$ where $\boldsymbol{\Omega}$ is assumed to meet the Gauss-Markov assumptions.

Given the parametric nature of the CNR, we can use likelihood methods\index{likelihood methods} for estimating the parameters. The likelihood function for each observation is

\begin{eqnarray}
\mathcal{L}_{i}  = f(y_{i}) = \frac{1}{\sqrt{2 \pi \sigma^{2}}}e^{-\frac{(y_{i} - \boldsymbol{x}'_{i}\boldsymbol{\beta})^{2}}{2\sigma^{2}}}
\end{eqnarray}

The maximum likelihood estimator\index{maximum likelihood estimator (MLE)|see{likelihood methods}} (MLE; $\hat{\boldsymbol{\beta}}_{mle}, \sigma^{2}_{mle}$) maximises $\log \mathcal{L}(\boldsymbol{\beta}, \sigma^{2})$. Since this is a function of $\boldsymbol{\beta}$ only through the sum of squared errors $SSE_{n}(\boldsymbol{\beta})$, maximising the likelihood is equivalent to minimising the sum of squared errors. Therefore

\begin{eqnarray}
\hat{\boldsymbol{\beta}}_{mle} = \boldsymbol{\hat{\beta}}_{ols}.
\end{eqnarray}


\begin{subappendices}
\section{Code for Spatial Stats in R}
\lstinputlisting{"Programs/UP504_OLS.R"}

\renewcommand\bibname{Further reading}

\bibliographystyle{econometrica}
\bibliography{Literature/OverLord}
\IfFileExists{Literature/ClassNormReg.bbl}{\chapter[The CNR model]{The classical normal regression (CNR) model}\label{chap_ClassNormReg}

\section{The CNR model}
\subsection{Basic set-up}
\subsection{The model in matrix notation}

\section{Different approaches to OLS}
\subsection{Least squares estimator}

\subsection{Geometric approach to OLS}
\subsubsection{Projection matrices and their relationships}
Let the matrix $\boldsymbol{P_{X}}$ be defined as

\begin{eqnarray}
\boldsymbol{P_{X}} \equiv \boldsymbol{X}(\boldsymbol{X}'\boldsymbol{X})^{-1}\boldsymbol{X}'.\label{eqn_CNR_Px}
\end{eqnarray}

$\boldsymbol{P_{X}}$ belongs to the general class of so-called projection matrices which have a number of interesting properties. Most importantly, the projection matrix $\boldsymbol{P_{X}}$ is a square matrix of dimension $n \times n$ that gives a vector space projection from $\mathbb{R}^{n}$ to a given subspace. Projection matrices are said to be \emph{idempotent} which implies that $\boldsymbol{P_{X}}\boldsymbol{P_{X}}=\boldsymbol{P_{X}}$.

This important property is straightforward to demonstrate in the case of equation \eqref{eqn_CNR_Px}, where we can use the fact that $$(\boldsymbol{X}'\boldsymbol{X})^{-1}\boldsymbol{X}'\boldsymbol{X}=\boldsymbol{I}$$ in order to show that $$\boldsymbol{P_{X}}\boldsymbol{P_{X}}=\boldsymbol{X}(\boldsymbol{X}'\boldsymbol{X})^{-1}\boldsymbol{X}'\boldsymbol{X}(\boldsymbol{X}'
\boldsymbol{X})^{-1}\boldsymbol{X}'=\boldsymbol{X}(\boldsymbol{X}'\boldsymbol{X})^{-1}\boldsymbol{X}'=\boldsymbol{P_{X}}$$. Furthermore, projection matrices can generally be classified either as \emph{orthogonal} projection matrices or \emph{oblique} (non-orthogonal) projection matrices.

In addition to being square, a projection matrix is also a symmetric matrix if and only if the vector space projection is orthogonal. $\boldsymbol{P_{X}}$ is an orthogonal projection matrix that spans the space $\mathcal{S}(\boldsymbol{X})$; consequently, any point in $\boldsymbol{X}$ that is projected by $\boldsymbol{P_{X}}$ will be projected onto itself. This implies that $\boldsymbol{P_{X}}\boldsymbol{X} = \boldsymbol{X}$. Again, this is easily verified by using equation \eqref{eqn_CNR_Px} which yields $\boldsymbol{P_{X}}\boldsymbol{X}=\boldsymbol{X}(\boldsymbol{X}'\boldsymbol{X})^{-1}\boldsymbol{X}'\boldsymbol{X} = \boldsymbol{X}$.

$\boldsymbol{M_{X}}$ is the orthogonal complement to $\boldsymbol{P_{X}}$, such that $\boldsymbol{M_{X}} = (\boldsymbol{I} - \boldsymbol{P_{X}})$. Consequently,

\begin{eqnarray}
\boldsymbol{P_{X}}\boldsymbol{M_{X}} &=& \boldsymbol{P_{X}}(\boldsymbol{I} - \boldsymbol{P_{X}}) \nonumber \\
&=& \boldsymbol{P_{X}}-\boldsymbol{P_{X}}\boldsymbol{P_{X}} = \boldsymbol{P_{X}}-\boldsymbol{P_{X}}=\boldsymbol{0}.
\end{eqnarray}

Since $\boldsymbol{P_{X}}$ and $\boldsymbol{M_{X}}$ are orthogonal to each other, any vector in $\mathcal{S}(\boldsymbol{X})$ will either be projected onto itself (i.e. $\boldsymbol{P_{X}}\boldsymbol{X} = \boldsymbol{X}$) or it will be annihilated ($\boldsymbol{M_{X}}\boldsymbol{X} = \boldsymbol{0}$).

Now, suppose that the matrix $\boldsymbol{X}$ can be decomposed into the submatrices $\boldsymbol{X}_{1}$ and $\boldsymbol{X}_{2}$, such that $\boldsymbol{X} = [\boldsymbol{X}_{1}\hspace{1ex}\boldsymbol{X}_{2}]$. Analogous to $\boldsymbol{P_{X}}$, we can now define the projection matrix $\boldsymbol{P}_{1} = \boldsymbol{X}_{1}(\boldsymbol{X}_{1}'\boldsymbol{X}_{1})^{-1}\boldsymbol{X}_{1}'$. This yields the following important relationships:

\begin{eqnarray}
\boldsymbol{P_{X}}\boldsymbol{P}_{1}=\boldsymbol{P}_{1} \qquad &\text{since}& \qquad \boldsymbol{P}_{1}<\boldsymbol{P_{X}}\\
\boldsymbol{M_{X}}\boldsymbol{M}_{1}=\boldsymbol{M_{X}} \qquad &\text{since}& \qquad \boldsymbol{M}_{1}>\boldsymbol{M_{X}}
\end{eqnarray}


\subsubsection{Geometry of the Frisch-Waugh-Lovell theorem}
See \citet{Lovell2008} for a simple proof that does not rely on matrix algebra. More comprehensive treatment of the Frisch-Waugh-Lovell theorem (FWL) can be found in \citet{Davidson2004}.

\subsection{General Method of Moments (GMM)}

\subsection{Maximum Likelihood Estimation and OLS}
The maximum likelihood estimate (MLE) is that assumed value of an unknown population parameter which maximises the probability of obtaining the sample data that is actually observed.

The normal regression model is the linear regression model under the restriction that the error term $e_{i}$ is independent of $\boldsymbol{x}_{i}$ and has the distribution $(0, \sigma^{2})$ Thus the standard set-up of OLS that we have just seen can easily be re-written as

\begin{eqnarray}
e_{i} | \boldsymbol{x}_{i} \sim (0, \sigma^{2})
\end{eqnarray}

Recall from section~\ref{chap_Review:Stats} that the mean and variance of a variable $V = \mu_{0} + e$ where $e \sim (0,\sigma^{2})$, is $\mu_{V} = \mu$ and $\sigma^{2}_{V}=\sigma^{2}_{e}$. This implies that

\begin{eqnarray}
y_{i} | \boldsymbol{x}_{i} \sim (\boldsymbol{x}_{i}'\boldsymbol{\beta}, \sigma^{2})
\end{eqnarray}

Equivalently, we can rewrite the linear model as $\boldsymbol{y} \sim N{\boldsymbol{\Omega}, \boldsymbol{\Omega}}$ where $\boldsymbol{\Omega}$ is assumed to meet the Gauss-Markov assumptions.

Given the parametric nature of the CNR, we can use likelihood methods\index{likelihood methods} for estimating the parameters. The likelihood function for each observation is

\begin{eqnarray}
\mathcal{L}_{i}  = f(y_{i}) = \frac{1}{\sqrt{2 \pi \sigma^{2}}}e^{-\frac{(y_{i} - \boldsymbol{x}'_{i}\boldsymbol{\beta})^{2}}{2\sigma^{2}}}
\end{eqnarray}

The maximum likelihood estimator\index{maximum likelihood estimator (MLE)|see{likelihood methods}} (MLE; $\hat{\boldsymbol{\beta}}_{mle}, \sigma^{2}_{mle}$) maximises $\log \mathcal{L}(\boldsymbol{\beta}, \sigma^{2})$. Since this is a function of $\boldsymbol{\beta}$ only through the sum of squared errors $SSE_{n}(\boldsymbol{\beta})$, maximising the likelihood is equivalent to minimising the sum of squared errors. Therefore

\begin{eqnarray}
\hat{\boldsymbol{\beta}}_{mle} = \boldsymbol{\hat{\beta}}_{ols}.
\end{eqnarray}


\begin{subappendices}
\section{Code for Spatial Stats in R}
\lstinputlisting{"Programs/UP504_OLS.R"}

\renewcommand\bibname{Further reading}

\bibliographystyle{econometrica}
\bibliography{Literature/OverLord}
\IfFileExists{Literature/ClassNormReg.bbl}{\input{Literature/ClassNormReg.bbl}}{\typeout{Need to run bibtex}}
\end{subappendices}
}{\typeout{Need to run bibtex}}
\end{subappendices}
}{\typeout{Need to run bibtex}}
\end{subappendices}
}{\typeout{Need bibtex on ClassNormReg}}
\end{subappendices}
